% ── §7c Methods: Worked Example ──
\section{Worked Example: Karnaugh Map Construction for PSICOV Target 1a3aA}
\label{sec:worked-example}

To make the pipeline concrete, we walk through the complete construction of a
Karnaugh map for positions 37 and 43 of PSICOV target 1a3aA (a 145-residue
protein with 5{,}378 alignment sequences).

\subsection{Step 1: Consensus Residues}
From the alignment, the majority residue at column~37 is glutamic acid (E,
He code~8), and at column~43 it is lysine (K, He code~13).  Both columns are
variable ($H(37) = 0.81$, $H(43) = 0.74$).

\subsection{Step 2: Gray-Code Transformation}
Applying the reflected binary Gray transformation $g(n) = n \oplus (n \gg 1)$:
\begin{align}
	g(8)  & = 8 \oplus 4 = 12 = 01100_2,  \\
	g(13) & = 13 \oplus 6 = 11 = 01011_2.
\end{align}
These five-bit codes determine the row (12) and column (11) of the primary
cell on the $32 \times 32$ Karnaugh map.  The corresponding flat cell index is
$c = 32 \times 12 + 11 = 395$.

\subsection{Step 3: Ternary Classification}
Scanning all 5{,}378 sequences, we find that the pair (E, K) at columns
(37, 43) is the reference (majority) combination, appearing in 1{,}834 of
5{,}378 sequences.  The cell at index $(12, 11)$ is therefore assigned
don't-care status ($-1$).  Other observed combinations---for example (D, R)
at cell index determined by $g(9) = 14$ for D and $g(14) = 9$ for R---are
assigned ON status ($+1$), while unobserved combinations receive OFF status
($0$).

\subsection{Step 4: Padding}
The $20 \times 20$ active region is embedded in a $32 \times 32$ grid by
assigning don't-care status to all cells with row or column index $\geq 20$.
After flattening, this produces a 1024-element ternary vector suitable for
Quine--McCluskey input.

\subsection{Step 5: Minimisation}
The popcount-grouped Quine--McCluskey algorithm processes the 1024-element
truth table, merging adjacent minterms and don't-cares to produce a set of
prime implicants.  For this particular pair, the minimiser identifies several
implicants including rules consistent with acidic--basic salt-bridge formation
at short-range separation.

\subsection{Step 6: Validation}
Positions 37 and 43 in 1a3aA are separated by only six residues in sequence,
placing them in separation bin~0 ($[6, 11)$).  Their C$\beta$ distance in the
native structure is 6.7~\AA, confirming that this pair is indeed a native
contact, which is consistent with its inclusion among the top-ranked pairs by
both our circuit scores and published DCA predictions.

\subsection{Interpretation}
This worked example illustrates how the entire analytical chain operates on
real data: consensus residues are identified from the MSA, transformed into
Gray-coded Karnaugh-map coordinates, classified as mutation/reference/unobserved,
padded to power-of-four dimensions, and minimised to produce interpretable
Boolean rules.  Each step is verified by unit tests and formal proofs,
ensuring that the output faithfully represents the input data under the
documented semantics.
